\chapter{六种实验条件参考}
\label{app:features}

\small
\begin{longtable}{p{0.24\textwidth}p{0.30\textwidth}p{0.31\textwidth}}
\caption{论文使用的六种显示实验条件。}\\
\toprule 条件 & 开关与输入 & 实际显示效果\\ \midrule\endfirsthead
\toprule 条件 & 开关与输入 & 实际显示效果\\ \midrule\endhead
Baseline & 关闭实验优化，展开全部节点，清空搜索和聚焦，缩放 100\%。 & 完整卡片和全部分支。\\
Compact Cards & 只开启 Compact。 & 全部节点保留，卡片只显示类型和短标题。\\
Optimized Overview & 开启 Compact 与 Semantic Zoom，缩放 50\%。 & 全部节点保留，使用最低细节查看整体结构。\\
Optimized Search & 在 Overview 上搜索固定唯一标题。 & 目标进入视口并保持显著，其他卡片淡化。\\
Subtree Focus & 开启 Compact，设置固定 Focus Root。 & 只显示指定分支及必要上下文。\\
Context Collapse & 开启 Compact，折叠固定目标路径之外的指定分支。 & 保留折叠根与摘要，隐藏其后代卡片。\\
\bottomrule
\end{longtable}

六种条件的详细步骤和测量目的见表\ref{tab:displayconditions}。插件中的其他显示功能不属于本论文实验条件。

\chapter{复现与证据索引}
\label{app:evidence}

以下路径均相对于仓库根目录：
\begin{description}
    \item[活动项目：] \path{testgame/testgame/}
    \item[分发插件：] \path{visual_scripting/addons/behavior_tree_editor/}
    \item[复杂树：] \path{testgame/testgame/behavior_trees/complex_guard_validation_tree.tres}
    \item[测试脚本：] \path{testgame/testgame/tests/}
    \item[多规模显示数据：] \path{testgame/testgame/test_results/multiscale_display_raw.csv}
    \item[多规模拖拽数据：] \path{testgame/testgame/test_results/multiscale_drag_raw.csv}
    \item[多规模汇总：] \path{testgame/testgame/test_results/multiscale_display_summary.csv} 与 \path{multiscale_drag_summary.csv}
    \item[分析与报告：] \path{research/display_optimization/analyze_multiscale_experiment.py} 与 \path{multiscale_experiment_results.md}
    \item[视觉截图：] \path{testgame/testgame/test_results/visual/}
    \item[可玩树原始数据：] \path{testgame/testgame/test_results/complex_display_experiment_raw.csv}
    \item[物理屏幕证据：] 实验根目录为 \path{research/display_optimization/}\ \path{physical_screen_size_experiment/}；其中 \path{data/2026-08-23_three_displays/} Session 包含 \path{combined_raw.csv}、\path{screen_size_comparison.csv} 与 \path{report_zh.md}。
    \item[发布包：] \path{dist/behavior-tree-editor-0.9.0.zip}
\end{description}

真实视觉回归必须使用 \texttt{--rendering-method gl\_compatibility}，不能使用 Dummy Headless，因为 Dummy 无法提供所需 SubViewport 纹理。

\chapter{初稿完成清单}

正式提交前，作者应：替换身份与导师占位符；按学校政策修改声明并如实说明 AI 协助；在 Tabula 核实截止时间和词数规则；请导师审阅文献、研究问题和证据边界；按院系要求执行 Turnitin；用模板编译并逐页检查；加入最终词数和要求的仓库/归档信息。
